\documentclass[10pt,a4paper]{amsart}
\usepackage[margin=19mm]{geometry}
\usepackage{amsmath,amssymb,mathtools}
\usepackage[scr=rsfs,cal=euler]{mathalfa}
\usepackage{xfrac}
\usepackage[unicode,hidelinks,bookmarks=false]{hyperref}
\newcommand{\ie}{i.e.\ }
\newcommand{\cf}{cf.\ }
\newcommand{\resp}{respectively\ }
\newcommand{\Subsection}{\subsection}
\providecommand{\nobreakdash}{\nobreak\hbox{-}}
\setlength{\emergencystretch}{2em}
\multlinegap0pt
\usepackage{mathtools}
\newtheorem{theorem}{Theorem}
\newtheorem{proposition}{Proposition}
\newcommand{\suppsec}[1]{the full-proof block at the end of this appendix}
\title{The Endogeneity of Miscalibration: \\ Impossibility and Escape in Scored Reporting}
\author{Lauri Lov\'{e}n, Sasu Tarkoma}
\date{Source: arXiv:2605.07671v2 (CC BY 4.0)}
\begin{document}
\maketitle

\noindent\emph{Source and licence.} The Endogeneity of Miscalibration: \\ Impossibility and Escape in Scored Reporting. Version arXiv:2605.07671v2 (CC BY 4.0), distributed by arXiv under the Creative Commons Attribution 4.0 International licence, http://creativecommons.org/licenses/by/4.0/, as recorded in the arXiv licence metadata for this exact version and shown on the abstract page https://arxiv.org/abs/2605.07671v2. Full licence notice: \texttt{licenses/arxiv-2605.07671-CC-BY-4.0.txt}.\par\smallskip

\noindent\textit{Editorial scope.} Counted unit: the article's theorem thm:reserve (source0.tex lines 397--403). The article's complete original proof of it is delivered with it (source0.tex lines 679--681, one \texttt{proof} environment of 3 lines, which the article places away from the statement). No mathematics is written, repaired or re-derived: the statement and the proof are the article's own lines, in the article's own notation, and no retained line is substituted or re-flowed.  Retained uncounted as the source-local context the proof needs: lines 431--443; lines 501--507.  Nothing was omitted from any retained span, so the package carries no omission record.  No retained line contains a citation, so the wrapper carries no bibliography and no bibliography entry.  No retained line contains a figure environment, an image-inclusion command or drawing code, so no image is delivered and the package declares no asset dependency.  Editorial formatting only, in addition: an AMS \texttt{amsart} preamble that restates the article's own declarations and macros used by the retained lines, namely the environments \texttt{theorem}, \texttt{proposition} on separate counters, together with the standard packages those lines need.  The retained environments print in this document as Theorem 1, Proposition 1 in order of appearance, whereas the article prints the same items with the numbering of its own full document.  The complete native source of the article as distributed in the arXiv e-print for this exact version is bundled unchanged as \texttt{sources/200106/source0.tex}.
\begin{theorem}[Cap, pinch, chord, and the critical slope]\label{thm:critical-slope}
Parts~(i)--(iii) concern any measurable $q : [0,1] \to [0,1]$ that admits a global best response at every type, as the conventions above require, and that attains first-best screening; part~(iv) quantifies over the slope-capped classes.
\begin{enumerate}
\item[(i)] \emph{(The cap.)} Assume (A1). Then $\gamma\, q(r) \le \beta\, D_G(p_{\min},r)$ for every $r \in [0,1]$, and moreover $q \equiv 0$ on $[0,p_{\min}]$; hence $q \le \mathrm{ramp}$ on all of $[0,1]$. First-best screening also forces weak feasibility, so $r_0$ is well defined, and then $q(r) < 1$ for every $r \in [p_{\min},r_0)$. The cap and the vanishing below $p_{\min}$ use only the weaker hypothesis that $\tau(p) = 0$ for $F$-almost every $p < p_{\min}$, which is the form the saturated regime of Proposition~\ref{prop:lambda-sat} needs.
\item[(ii)] \emph{(The pinch.)} Assume (A1) and strict feasibility. Then $\inf\{r : q(r) = 1\} = r_0$; the value $1$ is attained at $r_0$ itself or at points accumulating at $r_0$ from the right, strict right-accumulation being forced exactly when $q(r_0) < 1$; and if $q$ is continuous then $q(r_0) = 1$ exactly. No assumption is made that $\{q = 1\}$ is an interval.
\item[(iii)] \emph{(The chord bound.)} Assume (A2), strict feasibility, and $q$ continuous. Then for every $a \in [p_{\min},r_0)$,
\[
\frac{q(r_0)-q(a)}{r_0-a} \;\ge\; \frac{1-\mathrm{ramp}(a)}{r_0-a}, \qquad\text{hence}\qquad \mathrm{Lip}(q) \;\ge\; \Lambda_c(G) \;\ge\; \Lambda^*(G) \;>\; 0 .
\]
No differentiability of $q$ is used. If $q$ happens to be differentiable at $r_0$ from the left, then additionally $q'(r_0^-) \ge \Lambda^*(G)$.
\item[(iv)] \emph{(The threshold, both directions.)} Assume (A2) and strict feasibility. If $q \in \mathcal L_\Lambda$ attains first-best screening, then $\Lambda \ge \Lambda_c(G)$; continuity is automatic in $\mathcal L_\Lambda$. Conversely, if $\Lambda \ge \Lambda_c(G)$ then the truncated line $q_\Lambda(r) \coloneqq \operatorname{med}\{0,\ 1-\Lambda(r_0-r),\ 1\}$ belongs to $\mathcal Q_\Lambda$ and attains first-best screening, with a unique best response for every $p \ne p_{\min}$. Consequently $\delta_\Lambda(G) = \delta^{\mathcal L}_\Lambda(G) = 0$ \textbf{if and only if} $\Lambda \ge \Lambda_c(G)$, and $\Lambda_c$ is attained, so the infimum over admissible slopes is a minimum.
\end{enumerate}
\end{theorem}

\begin{proposition}[Exact screening at reduced intensity]\label{prop:lambda-sat}
Assume (A1) and $\gamma/\beta > D_G(p_{\min},1)$. Put $\lambda_{\mathrm{sat}} \coloneqq \beta D_G(p_{\min},1)/\gamma \in (0,1)$ and $q_{\mathrm{sat}} \coloneqq \lambda_{\mathrm{sat}}\,\mathbf 1\{r \ge 1\}$. Then $q_{\mathrm{sat}}$ induces $\tau(p) = \lambda_{\mathrm{sat}}\,\mathbf 1\{p \ge p_{\min}\}$ for every $p \ne p_{\min}$, so the screening boundary lands exactly at the first-best cutoff, and
\[
W(q_{\mathrm{sat}}) \;=\; u_d + \lambda_{\mathrm{sat}}\int_{p_{\min}}^1 \Pi(p) f(p)\, dp \;>\; u_d .
\]
This is optimal within the class of approval rules that reject almost every type below $p_{\min}$.
\end{proposition}

\begin{theorem}[Hard oversight at $r_0(G)$]\label{thm:reserve}
Assume (A1) and weak feasibility $\gamma/\beta \le D_G(p_{\min},1)$, and let $r_0 \in (p_{\min},1]$ be the unique solution of $D_G(p_{\min},r_0) = \gamma/\beta$. Then:
\begin{enumerate}
\item[(i)] The step rule $q(r) = \mathbf 1\{r \ge r_0\}$ induces $\tau(p) = \mathbf 1\{p \ge p_{\min}\}$ for every $p \ne p_{\min}$, with strict optimality of the induced report; at $p = p_{\min}$ the agent is exactly indifferent between $r = p_{\min}$ and $r = r_0$, a single ($F$-null) type. Hence the step rule attains first-best welfare under any tie-breaking.
\item[(ii)] The reserve is an affine function of $\sqrt{\gamma/\beta}$, that is $r_0 = p_{\min} + \kappa_0\sqrt{\gamma/\beta}$ for a constant $\kappa_0$ independent of $(p_{\min},\gamma/\beta)$ across the whole feasible range, if and only if $G$ is quadratic, and $\kappa_0 = 1$ if and only if $G'' \equiv 2$, the Brier normalization.
\end{enumerate}
\end{theorem}

\begin{proof}[Proof sketch]
The cap of Theorem~\ref{thm:critical-slope}(i) forces an affine rule to vanish on $[0,p_{\min}]$, hence to vanish identically, hence to approve nobody, which is the regime-free screening half. Under weak feasibility the step rule of Theorem~\ref{thm:reserve}(i) attains $W^*$ while $\tau\Pi\leq\max(\Pi,0)$ caps every rule at $W^*$, so an optimum exists, attains first-best screening and is therefore not affine. The full verification is in \suppsec{S.4}.
\end{proof}

\end{document}
